\section{前言}

我14岁的时候，正巧赶上1997年大牛市的末端。或许是由于身边亲戚赚钱案例的吸引，或许是受到股市类题材电视剧的启发，总之，我这个从来都对自己的智商不缺乏自信的小毛孩子，幻想着可以从股票市场上赚到不小的财富。正是由于我的年少轻狂，才将我的家庭引入了深渊。记得那时开立一个股票账户至少需要8万元的保证金，而我也只能鼓动我老爸来开立股票账户。由于在大牛市的顶部入市，我家买入的股票立刻就遭到了套牢，之后为了实现“解套”而不得不投入越来越多的资金。很显然，更多的投入只是带来了更多的亏损，最终我老爸实在没有勇气继续面对这个可怕的市场了。这样一来，我家的财政大权就完全落在了我的手中，而我的接手同样无法扭转继续亏损的轨迹。通过多年的“努力”，在我19岁的时候就已经将爸妈毕生的积蓄亏损过半。

股票市场的确伴随着我的成长，也让我丧失了很多同龄人的乐趣。好的一面是，在股票市场上碰的钉子让我成熟了不少，甚至培养起了更为可贵的责任心。或许是因为年轻人的思维更加灵活，而不会像中年人那样有石膏般僵硬的思维误区，最终我还是一步一步走出了泥潭，只不过这花费了我8年的光阴。

多年以后的今天，我虽然赚到了数十倍的利润，甚至可以稳定地在市场上赚取利润，但我仍旧在反思，走交易这条路真的值得吗？曾经在市场上的痛苦挣扎，让我至今都心有余悸。很显然，我绝不可能具有金融从业者那种必须具备的（对交易的）乐观态度。我不是金融行业的从业人员，只是为自己负责地管理着一笔自有资金，甚至我最终还找了一份与金融行业完全无关的工作。在交易的成长之路上，我受到很多交易书籍的宝贵启发，并在交易论坛的讨论中产生了大量的共鸣。我得到很多网友的认可，而大家的支持的确让我非常感动，这让我在这个冰冷的事业中感到了可贵的温暖。在很多热心网友的支持与帮助下，我在网络上创建了一个非营利性的论坛——“交易者之家”。道理很简单，如果你把自己的博弈对手当做自己的老师，那么你几乎毫无取胜的可能，这无异于奢望从广告中学会消费。我希望“交易者之家”能够成为脱离行业内利益动机的交流净土，从而实现一种交易者之间单纯的互助模式。论坛上的每个人都在无私地分享自己的经验与感悟，随着不断地交流与总结，本书的雏形也渐渐地形成了。从某种意义上来讲，本书是我的无心之作。我的文章被众多网友自发地归纳、整理、转发，而更多人的期待让我有信心将我的交易思想系统地整理成书。

我觉得本书最大的看点在于，这是一本从交易者的视角观察市场的作品。或者说，这是一本真正意义上为交易者而写的书。事实上，这个行业内的每个人所阐述的观点，都不可能不受到自身角色的影响。很显然，我与读者朋友们没有丝毫的不同，都只不过是市场中一位渺小的交易者。或许，我比那些高不可攀的业内人士更加了解中小投资者所处的自身环境，以及思想中的误区。最低限度，我不希望再为读者朋友们增添任何新的误区。本书并不想证明什么，只是想告诉大家一些市场的真相。

我相信，市场上已经不可能再存在任何新鲜的东西了。但我仍希望能够通过新鲜的视角来思考市场中的老问题，否则本书也就不再具有任何存在的意义了。我毕竟不是职业作家，那么我就很难给大家优雅的词藻，以及严谨的逻辑性。甚至由于本书的创作经历了太长的时间跨度，从而使得内容结构缺乏层次感，并且很多地方存在重复的嫌疑。但最低限度，我有义务让每一位读者感受到真诚与客观。我所谓的客观，并非是指观点上的绝对正确，恐怕在市场上没有什么结论是绝对正确的。相反，任何结论的有效性都需要建立在特定角度的基础上。我只不过是希望本人的观点，不会受到任何动机的影响而偏离了公正与客观。

每一本关于交易的书籍，读者都是最终的消费者。那么读者就成了作者的上帝，没有哪个作者会不顾上帝的感受。那么我是给大家需要的东西，还是给大家想要的东西呢？好在我既不想出名，也不想通过本书给自己带来任何经济上的利益。我可以很诚实地告诉大家，本书并不能通过什么独门秘籍，帮助你在明天的市场上多赚一些利润。本书并不会为了销量而迁就读者朋友们的本能需要。相反，本书的核心就是通过我曾经走过的弯路，来挖掘交易者思想上的误区。我不惧怕直视自己思想与心灵上脆弱的一面，我同样也希望读者朋友们能够像我一样坚强。
